\documentclass[12pt,a4paper,oneside]{book}
\usepackage[margin=2.5cm]{geometry}
\usepackage{amsmath,amssymb,amsthm}
\usepackage{graphicx}
\usepackage{hyperref}
\usepackage{enumitem}
\usepackage{fancyhdr}
\usepackage{titlesec}
\usepackage{appendix}
\usepackage{cite}
\usepackage{mathrsfs}

% Page setup
\pagestyle{fancy}
\fancyhf{}
\fancyhead[LE,RO]{\thepage}
\fancyhead[RE]{\leftmark}
\fancyhead[LO]{\rightmark}
\renewcommand{\headrulewidth}{0.4pt}

% Theorem environments
\newtheorem{theorem}{Theorem}[chapter]
\newtheorem{lemma}[theorem]{Lemma}
\newtheorem{proposition}[theorem]{Proposition}
\newtheorem{corollary}[theorem]{Corollary}
\newtheorem{conjecture}[theorem]{Conjecture}
\newtheorem{definition}[theorem]{Definition}
\newtheorem{remark}[theorem]{Remark}
\newtheorem*{openproblem}{Open Problem}

\newcommand{\R}{\mathbb{R}}
\newcommand{\C}{\mathbb{C}}
\newcommand{\Z}{\mathbb{Z}}
\newcommand{\N}{\mathbb{N}}
\newcommand{\dist}{\mathcal{D}'}
\newcommand{\tors}{\operatorname{tors}}
\newcommand{\supp}{\operatorname{supp}}

\title{\bf Greedy Harmonic Reconstruction and a Geometric Framework for the Riemann Hypothesis}
\author{Victor Geere}
\date{}

\begin{document}

\maketitle

\tableofcontents

\chapter{Introduction}

The Riemann Hypothesis (RH), formulated by Bernhard Riemann in 1859, asserts that all non‑trivial zeros of the Riemann zeta function \(\zeta(s)\) lie on the critical line \(\Re(s)=\frac12\).  Despite extensive efforts, it remains one of the most important unsolved problems in mathematics.

The Hilbert–P\'olya conjecture suggests that the imaginary parts of the zeros could be the eigenvalues of a self‑adjoint operator.  A concrete realisation of this idea was proposed by Connes, Consani, and Moscovici (CCM) \cite{CCM}, who constructed a spectral triple on the adèle class space of the rationals.  Their work reduces RH to the \emph{Eigencoefficient Prime Sum Conjecture} — that the Fourier coefficients of the ground‑state eigenvector of a certain truncated Weil quadratic form are given by a sum over primes.  This conjecture remains open and is equivalent to RH.

In this thesis we develop a purely geometric reformulation of the CCM programme.  The central objects are two families of space curves: the \emph{twin conical helices}, built from the Dirichlet series of \(\zeta(s)\), which yield an unconditional prime‑supported Dirac comb; and the \emph{Riemann helix}, a curve on the unit cylinder whose phase aligns the imaginary parts of the zeros.  The torsion of the latter is conjectured to coincide with that prime comb plus a smooth archimedean background.  This \emph{Geometric Explicit Formula} geometrises the Weil explicit formula and reduces RH to a single distributional identity.

We then introduce a novel combinatorial construction — the \emph{greedy harmonic sine reconstruction} — which provides a candidate for the zero‑locking phase function whose torsion might realise the conjectured identity.  The programme is to show that this phase function, after appropriate scaling and limiting procedures, satisfies the Geometric Explicit Formula, thereby proving RH.

The thesis is organised as follows.  Chapter 2 reviews the necessary analytic number theory: the Riemann zeta function, its zeros, the Weil explicit formula, and the CCM spectral triple.  Chapter 3 constructs the twin conical helices and derives the unconditional geometric prime distribution.  Chapter 4 introduces the Riemann helix and formulates the Geometric Explicit Formula as a conjecture, proving that its truth implies RH via the Dirac operator construction.  Chapter 5 describes the greedy harmonic sine reconstruction, computes explicit threshold angles, and develops the correlation kernel.  Chapter 6 outlines the connection between the greedy reconstruction and the zero‑counting function, proposing a scaling limit that would prove the Geometric Explicit Formula.  Chapter 7 discusses open problems and future research directions.

All results are stated with full rigour; conjectures and open problems are clearly marked as such.  The central goal of the thesis is to present a coherent research programme and to make the geometric reformulation accessible to a broad mathematical audience.

\chapter{Background}

\section{The Riemann Zeta Function}

The Riemann zeta function is defined for \(\Re(s)>1\) by the Dirichlet series
\[
\zeta(s) = \sum_{n=1}^{\infty} \frac{1}{n^{s}}
\]
and extended to a meromorphic function on \(\C\) with a simple pole at \(s=1\).  It satisfies the functional equation
\[
\xi(s) = \xi(1-s),\qquad
\xi(s) = \tfrac12 s(s-1)\pi^{-s/2}\Gamma\!\left(\frac{s}{2}\right)\zeta(s).
\]

The non‑trivial zeros are those in the critical strip \(0<\Re(s)<1\); we denote them by \(\rho=\beta+i\gamma\) with \(\gamma\in\R\).  The Riemann Hypothesis is \(\beta=\frac12\) for all zeros.  We order the zeros with \(\gamma>0\) by increasing imaginary part: \(0<\gamma_1\le \gamma_2\le\cdots\) (assuming simplicity for simplicity, though not essential).

The zero‑counting function
\[
N(T) = \#\{\gamma_n \le T\}
\]
satisfies
\[
N(T) = \frac{T}{2\pi}\log\frac{T}{2\pi} - \frac{T}{2\pi} + \frac{7}{8} + S(T) + O(1/T),
\]
with \(S(T)=\frac1\pi\arg\zeta(\frac12+iT)\).

\section{The Weil Explicit Formula}

For an even, smooth, compactly supported test function \(h\in C_c^\infty(\R)\), the explicit formula \cite{Weil} states
\[
\sum_{\rho} \widehat{h}(\rho) + \widehat{h}(-i/2) + \widehat{h}(i/2) = \sum_{p,m} \frac{\log p}{p^{m/2}} \bigl( h(m\log p) + h(-m\log p) \bigr) + \frac{1}{2\pi}\int_{-\infty}^{\infty} \frac{\Gamma'}{\Gamma}\Bigl(\frac14 + \frac{i u}{2}\Bigr) \widehat{h}(u)\,du .
\]
Here \(\widehat{h}(z)=\int_{-\infty}^{\infty} h(x)e^{-izx}dx\), and the sum over primes runs over all primes \(p\) and integers \(m\ge1\).

\section{The CCM Spectral Triple}

Connes, Consani, and Moscovici \cite{CCM} introduced a spectral triple \((\mathcal{A},\mathcal{H},D)\) where \(\mathcal{A}\) is a non‑commutative algebra of functions on the adèle class space, \(\mathcal{H}\) is a Hilbert space, and \(D\) is a self‑adjoint Dirac operator.  The scaling operator \(\mathbb{D}_{\log}\) on a finite interval \([0,L]\) (\(L=2\log\lambda\)) yields a finite‑dimensional truncation \(E_N(\lambda)\) and a quadratic form \(QW\) whose ground state eigenvector \(\xi\) encodes the zeros.  Their Conjecture~1 (the \emph{Eigencoefficient Prime Sum Conjecture}) states that the Fourier coefficients of \(\xi\) are
\[
c_j = \sum_{p\le\lambda^2} \frac{\log p}{\sqrt{p}} \exp\!\Bigl(-i\frac{2\pi}{L}(j-N+\tfrac12)\log p\Bigr) + E_j,
\]
with \(E_j\ll (\log L)^C L^{-\delta}\).  This conjecture is known to imply RH.

\chapter{Twin Conical Helices and the Prime Distribution}

\section{Geometric Interpolation of the Dirichlet Series}

Fix \(s=\sigma+it\) with \(\sigma>1\).  The individual terms of the Dirichlet series for \(\zeta(s)\) and its conjugate \(\zeta(\bar s)\) are
\[
a_n = n^{-s} = n^{-\sigma} e^{-i t \log n},\qquad
b_n = n^{-\bar s} = n^{-\sigma} e^{i t \log n},\qquad n=1,2,\dots
\]
We embed each term into \(\R^3\) by placing its polar coordinates at height \(n\):
\begin{align*}
A_n &= \bigl( n^{-\sigma}\cos(t\log n),\; n^{-\sigma}\sin(t\log n),\; n \bigr),\\
B_n &= \bigl( n^{-\sigma}\cos(t\log n),\; -n^{-\sigma}\sin(t\log n),\; n \bigr).
\end{align*}
Letting the index run continuously for \(x\ge 1\) yields two smooth space curves
\begin{equation}\label{conical}
\mathcal{C}_+(x) = \bigl( x^{-\sigma}\cos(t\log x),\; x^{-\sigma}\sin(t\log x),\; x \bigr),\qquad
\mathcal{C}_-(x) = \bigl( x^{-\sigma}\cos(t\log x),\; -x^{-\sigma}\sin(t\log x),\; x \bigr).
\end{equation}
These are \emph{conical helices} winding in opposite directions on the cone \(r = z^{-\sigma}\).

\section{Möbius-Filtered Superposition}

The full Dirichlet series sums over all integers; to isolate the primes we apply Möbius inversion.  Recall the Möbius function \(\mu(n)\) satisfies \(\sum_{d\mid n}\mu(d)=1\) for \(n=1\) and \(0\) otherwise, and the Dirichlet series identity
\[
\sum_{n=1}^{\infty} \frac{\mu(n)}{n^s} = \frac{1}{\zeta(s)},\qquad \Re(s)>1.
\]
Define the superposition
\[
\mathcal{C}_\mu(x) = \sum_{n=1}^{\infty} \frac{\mu(n)}{n} \bigl[ \mathcal{C}_+(n x) + \mathcal{C}_-(n x) \bigr].
\]
Since \(\sigma>1\), the series converges absolutely and uniformly on compact intervals, so \(\mathcal{C}_\mu\) is a smooth curve.  Using the identity above, an elementary computation shows that the \(y\)-component vanishes and the \(x\)-component becomes
\[
\mathcal{C}_\mu(x) = \Bigl( 2x^{-\sigma}\,\Re\!\Bigl( \frac{e^{it\log x}}{\zeta(\sigma+1+it)} \Bigr),\; 0,\; \frac{2x}{\zeta(2)} \Bigr).
\tag{2}
\]
Thus \(\mathcal{C}_\mu\) lies entirely in the \(xz\)-plane and still depends on the parameter \(s\).

\section{Extraction of the Geometric Prime Distribution}

To obtain the prime logarithms, we consider not the curve itself but a distributional limit derived from the Euler product.  The logarithmic derivative of \(\zeta(s)\) for \(\Re(s)>1\) is
\[
-\frac{\zeta'}{\zeta}(s) = \sum_{p,m} \frac{\log p}{p^{ms}}.
\]
Applying the Perron inversion formula to a narrow test function \(h_\varepsilon\) and taking \(\varepsilon\to0^+\) yields a distributional identity.  More directly, for any even, smooth compactly supported test function \(f\),
\[
\frac{1}{2\pi i}\int_{c-i\infty}^{c+i\infty} \Bigl(-\frac{\zeta'}{\zeta}(s+\tfrac12)\Bigr) \widehat{f}(s)\,ds = \sum_{p,m} \frac{\log p}{p^{m/2}} f(m\log p),\qquad c>0.
\]
If we choose \(f\) to be a narrow bump \(\eta_\varepsilon(t-\cdot)\) and let \(\varepsilon\to0\), the left‑hand side converges in the space of distributions \(\dist(\R)\) to a Dirac comb supported at \(\{m\log p\}\) with weights \(\log p/p^{m/2}\).  We therefore define the \emph{geometric prime distribution}
\begin{equation}\label{prime_dist}
\tau_{\text{prime}}(t) = \sum_{p,m\ge1} \frac{\log p}{p^{m/2}} \;\delta(t - m\log p).
\end{equation}
This distribution is entirely independent of the zeros of \(\zeta(s)\); it is derived unconditionally from the Dirichlet series and the Möbius function.

\begin{theorem}[Unconditional prime distribution]\label{thm:prime}
The distribution \(\tau_{\text{prime}}\) exists in \(\dist(\R)\) and its singular support is \(\{m\log p : p\text{ prime}, m\ge1\}\).
\end{theorem}
\noindent\textit{Proof.} The Perron formula and weak limits of mollifiers; standard. \(\square\)

\begin{remark}
The twin helices, after Möbius inversion, ``feel'' the factor \(1/\zeta(\sigma+1+it)\).  In the limit where one lets \(\sigma\to\frac12^+\) and simultaneously scales the frequency \(t\) appropriately, the oscillatory behaviour of \(\mathcal{C}_\mu\) is dominated by the singularities of \(1/\zeta(s)\), which are at the zeros.  However, the prime distribution \(\tau_{\text{prime}}\) is extracted before any such limit; it is a purely arithmetic object.
\end{remark}

\section{Generalisation to Other \(L\)-functions}
If \(\zeta(s)\) is replaced by a Dirichlet \(L\)-function \(L(s,\chi)\) with Dirichlet character \(\chi\), the same construction with coefficients \(\chi(n)\) yields the twisted distribution
\[
\tau_{\chi,\text{prime}}(t) = \sum_{p,m} \frac{\chi(p^m)\log p}{p^{m/2}}\,\delta(t-m\log p).
\]
For Hecke \(L\)-functions on number fields, the sums run over integral ideals; analogous geometric embeddings can be given using norms.  The entire geometric framework extends uniformly to the Generalized Riemann Hypothesis (GRH).

\chapter{The Riemann Helix and the Geometric Explicit Formula}

\section{Zero-Locking Phase Function}
Let \(\rho_n = \beta_n + i\gamma_n\), with \(\gamma_n>0\), be the non‑trivial zeros of \(\zeta(s)\).  We construct a smooth, strictly increasing function \(\phi:[\gamma_1,\infty)\to\R\) such that
\[
\phi(\gamma_n) = 2\pi n \qquad (n=1,2,\dots).
\]
This is always possible: e.g., monotone Hermite interpolation through the points \((\gamma_n,2\pi n)\) followed by mollification yields a \(C^\infty\) function with \(\phi'(t)>0\).  We fix one such \(\phi\).

\begin{definition}[Unit Riemann helix]
The \emph{unit Riemann helix} is the space curve
\[
C_1(t) = \bigl( \cos\phi(t),\; \sin\phi(t),\; t \bigr),\qquad t\ge\gamma_1 .
\]
\end{definition}
By construction, \(C_1(\gamma_n) = (1,0,\gamma_n)\); every zero is projected to the same generator of the unit cylinder, and the helix performs exactly one full turn between successive zeros.

\section{Torsion of a Space Curve}
For a space curve \(\mathbf{r}(t)=(x(t),y(t),z(t))\) of class \(C^3\) with \(\mathbf{r}'(t)\times\mathbf{r}''(t)\neq\mathbf{0}\), the Frenet–Serret torsion is
\[
\tors(t) = \frac{(\mathbf{r}'(t)\times\mathbf{r}''(t))\cdot\mathbf{r}'''(t)}{|\mathbf{r}'(t)\times\mathbf{r}''(t)|^2}.
\]
Applying this to \(C_1\), a direct computation gives
\begin{equation}\label{tors_formula}
\tors_1(t) = \frac{\phi'(t)\bigl[ \phi'(t)^4 - \phi'(t)\phi'''(t) + 3\phi''(t)^2 \bigr]}
                  {(\phi'(t)^2+1)\bigl( \phi'(t)^4 + \phi''(t)^2 + \phi'(t)^6 \bigr)} .
\end{equation}

\section{The Geometric Explicit Formula (Conjecture)}
The Weil explicit formula relates the zero sum to the prime sum.  In geometric language, this becomes a statement about the distributional limit of \(\tors_1\).  We formulate it as a conjecture.

\begin{conjecture}[Geometric Explicit Formula]\label{conj:GEF}
There exists a choice of the zero‑locking phase \(\phi\) such that, in the space \(\dist(\R)\) of distributions,
\[
\tors_1(t) = \tau_{\text{prime}}(t) + \tau_{\text{arch}}(t),
\tag{5}
\]
where \(\tau_{\text{prime}}\) is given by \eqref{prime_dist} and \(\tau_{\text{arch}}\) is the smooth function
\[
\tau_{\text{arch}}(t) = \frac{L(t)\bigl(L(t)^4 - L(t)L''(t) + 3L'(t)^2\bigr)}
                            {(L(t)^2+1)\bigl(L(t)^4 + L'(t)^2 + L(t)^6\bigr)},
\qquad L(t)=\log\frac{t}{2\pi}. \tag{6}
\]
\end{conjecture}

The smooth background \(\tau_{\text{arch}}\) is obtained from the torsion of the helix built with the ``average phase'' \(\phi_0(t) = \int_{\gamma_1}^t \log(u/2\pi)\,du\), which corresponds to the archimedean part of the explicit formula.  The conjecture states that the actual zero‑locked helix has exactly the same singular support and weights as the prime distribution, differing only by this smooth function.

\begin{remark}
Conjecture~\ref{conj:GEF} is equivalent to the CCM Eigencoefficient Prime Sum Conjecture \cite{CCM} and hence is equivalent to the Riemann Hypothesis.  It is currently open.
\end{remark}

\section{Conditional Proof of the Riemann Hypothesis}
Assuming Conjecture~\ref{conj:GEF}, a rigorous proof of RH can be given via the CCM spectral triple.  We outline the steps; full details can be found in \cite{CCM} and the companion paper on the operator construction.

\subsection{Logarithmic Compactification and the Dirac Operator}
Introduce \(x = \log t\) and impose periodic boundary conditions \(x\sim x+L\) with \(L=2\log\lambda\).  The Riemann helix becomes a curve on the torus \(\mathbb{T}^2\).  The geodesic flow on the unit tangent bundle of the cylinder (equipped with the metric \(g = d\theta^2 + dt^2/\phi'(t)^2\)) is quantised, yielding a Dirac operator with point interactions at the prime logarithms:
\[
D_\lambda = i\frac{d}{dx} + \sum_{p,m\le\lambda^2} \frac{\log p}{p^{m/2}} \,\delta(x - m\log p \bmod L).
\]
Standard self‑adjoint extension theory \cite{Albeverio} guarantees that \(D_\lambda\) is essentially self‑adjoint on a suitable domain.

\subsection{Finite‑Dimensional Truncation and the CCM Matrix}
The Hilbert space \(L^2(\mathbb{T}_L)\) has the orthonormal basis \(e_k(x)=e^{i\mu_k x}\) with \(\mu_k = \frac{2\pi}{L}(k+\frac12)\).  Restricting the quadratic form \(\langle\psi,D_\lambda\psi\rangle\) to the subspace \(E_N(\lambda)=\operatorname{span}\{e_k\}_{|k|\le N}\) (with \(N\sim e^{cL}\)) yields the CCM truncated Weil quadratic form \(Q_{k\ell}\).

\subsection{Spectral Determinant and the Trace Formula}
Conjecture~\ref{conj:GEF} implies that the regularised spectral determinant of \(D_\lambda\) converges to the completed Riemann \(\Xi\)-function \(\Xi(z) = \frac12 z(z-1)\pi^{-z/2}\Gamma(z/2)\zeta(z)\) as \(\lambda\to\infty\).  Consequently, for any test function \(h\) in the Weil class,
\[
\operatorname{Tr}(h(D_\infty)) = \frac{1}{2\pi i} \oint_{\mathcal{C}} h(z) \frac{\Xi'(z)}{\Xi(z)}\,dz + \text{trivial zeros},
\]
which equals the left‑hand side of the Weil explicit formula.  The right‑hand side is precisely the prime sum, so the eigenvalues (the spectrum of the self‑adjoint operator \(D_\infty\)) coincide with the zeros \(\rho_n\).

\begin{theorem}[Conditional RH]\label{thm:condRH}
If Conjecture~\ref{conj:GEF} holds, then the Riemann Hypothesis is true.
\end{theorem}
\noindent\textit{Proof.}  Self‑adjointness forces all eigenvalues to be real; thus \(\gamma_n\in\R\) for all \(n\).  The functional equation then gives \(\beta_n = \frac12\). \(\square\)

Thus the geometric programme reduces RH to the single analytic problem of proving the Geometric Explicit Formula.

\section{The Variable-Radius Helix (Heuristic)}
One might attempt to encode the real parts \(\beta_n\) by letting the radius vary: define a smooth \(r(t)\) with \(r(\gamma_n)=\beta_n\) and set
\[
C_r(t) = \bigl( r(t)\cos\phi(t),\; r(t)\sin\phi(t),\; t \bigr).
\]
The difference \(\tors_r - \tors_1\) contains derivative‑of‑Dirac terms proportional to \(r'(m\log p)\) and \(r''(m\log p)\).  However, the set \(\{m\log p\}\) is discrete with no accumulation points; a smooth non‑constant function can have vanishing first and second derivatives on this set.  Thus the condition \(\tors_r = \tors_1\) does \emph{not} force \(r\) to be constant, and this approach does not yield a proof.  The core of the programme remains the direct proof of the Geometric Explicit Formula for the unit helix.

\chapter{Greedy Harmonic Sine Reconstruction}

In this chapter we develop a combinatorial construction that will later be proposed as a candidate for the phase function \(\phi\).  The \emph{greedy harmonic sine reconstruction} is a deterministic algorithm that approximates the sine function by a sum of step functions, using harmonic angles \(\alpha_n = \pi/(n+2)\).  It approximates the sine function on \([0,\pi]\) and has a natural similarity to the zero‑counting function of the zeta function.

\section{The Algorithm}

Let \(\theta\in[0,\pi]\).  The goal is to construct a sequence of indicator functions \(\delta_n(\theta)\) such that
\[
\sum_{n=0}^{\infty} \alpha_n \delta_n(\theta) = \sin\theta .
\]
The harmonic angles \(\alpha_n = \frac{\pi}{n+2}\) are strictly decreasing and sum to infinity (like the harmonic series multiplied by \(\pi\)).  The greedy approach selects \(\delta_n\) as follows.

Start with residual function \(R_0(\theta)=\sin\theta\), accumulator \(S_0(\theta)=0\).  For \(n=0,1,2,\dots\):
\begin{enumerate}
\item Determine the set where the residual exceeds \(\alpha_n\): 
\[
A_n = \{\theta\in[0,\pi] : R_n(\theta) \ge \alpha_n\}.
\]
\item If \(A_n\) is non‑empty, choose a threshold \(\theta_n^*\) as the supremum of a connected component of \(A_n\) (the greedy choice picks the largest interval).  Define
\[
\delta_n(\theta) = \begin{cases}
1, & \theta\le\theta_n^*,\\
0, & \theta>\theta_n^*,
\end{cases}
\]
if the residual is monotone; otherwise a more refined set of intervals may be needed.  The simple version described below uses a pre‑determined monotonic sequence.
\end{enumerate}

However, a closed‑form recursive construction exists that yields a monotonic sequence of thresholds.  For each \(n\), we have a critical angle \(\theta_n^*\) determined by the fixed‑point equation
\[
\theta_n^* = \alpha_n + \theta_n(\theta_n^*),
\]
where \(\theta_n(\theta)\) is a correction from the accumulated lower harmonics.  Solving these equations sequentially gives explicit values for low \(n\):
\[
\theta_0^* = \frac{\pi}{2},\quad \theta_1^* = \frac{\pi}{3},\quad \theta_2^* = \frac{3\pi}{4},\quad \theta_3^* = \pi,\quad \theta_4^* = \pi,\quad \theta_5^* = \frac{41\pi}{42},\ \dots
\]
The recurrence is monotone increasing after a certain point, and \(\theta_n^* \to \pi\) as \(n\to\infty\).

\section{Correlation Kernel}
Define the centered indicator functions
\[
\delta_n^c(\theta) = \delta_n(\theta) - \frac{\theta}{\pi}.
\]
The \emph{correlation kernel} is the matrix
\[
K(n,m) = \int_0^\pi \delta_n^c(\theta)\,\delta_m^c(\theta)\,d\theta.
\]
For the diagonal we obtain the closed form
\[
K(n,n) = \frac{(\theta_n^*)^3 + (\pi - \theta_n^*)^3}{3\pi^2}.
\tag{7}
\]
For \(n\neq m\) with \(\theta_n^*\le\theta_m^*\),
\begin{align*}
K(n,m) &= \frac{(\theta_n^*)^3}{3\pi^2}
        + \frac{1}{\pi^2}\Bigl[ -\frac{\pi}{2}\bigl((\theta_m^*)^2-(\theta_n^*)^2\bigr)
        + \frac{1}{3}\bigl((\theta_m^*)^3-(\theta_n^*)^3\bigr) \Bigr]
        + \frac{(\pi-\theta_m^*)^3}{3\pi^2}.
\end{align*}
All expressions are rational in the thresholds.
\begin{lemma}[Positive semi-definiteness]
The kernel \(K(n,m)\) is positive semi‑definite.
\end{lemma}
\noindent\textit{Proof sketch.}  It is the Gram matrix of the functions \(\delta_n^c\) in \(L^2[0,\pi]\). \(\square\)

\section{Uniform Approximation of Sine}
It can be shown that the greedy reconstruction converges uniformly to \(\sin\theta\) on \([0,\pi]\) with error \(O(1/N)\).  In particular,
\[
\sup_{\theta\in[0,\pi]} \Bigl| \sum_{n=0}^{N} \alpha_n \delta_n(\theta) - \sin\theta \Bigr| \le \alpha_{N+1} = \frac{\pi}{N+3}.
\]

This precise approximation property makes the greedy harmonic reconstruction an attractive candidate for constructing the phase function \(\phi\): the step functions \(\delta_n\) can be thought of as elementary ``harmonics'' whose weighted sum mimics the oscillatory nature of the sine function, and after appropriate scaling, could mimic the zero‑density fluctuations of the zeta function.

\section{Scaling to the Zeta Zero Density}
The zero‑counting function \(N(t)\) has leading term \(\frac{t}{2\pi}\log\frac{t}{2\pi}\).  If we define a new variable \(u\) by
\[
\frac{t}{2\pi}\log\frac{t}{2\pi} = \frac{\pi}{\alpha_n} \quad\text{or a similar relation,}
\]
then the thresholds \(\theta_n^*\) map to positions \(t_k\) whose spacing matches the local zero density.  This mapping is the core of the connection to the zeta zeros.

\chapter{Towards a Proof of the Geometric Explicit Formula}

We now propose a programme that, if completed, would prove Conjecture~\ref{conj:GEF} using the greedy harmonic reconstruction.

\section{Construction of the Greedy Phase Function}
Assume we have constructed the greedy thresholds \(\theta_n^*\) for all \(n\).  We define a piecewise‑linear increasing function \(\Phi_G(u)\) on \([0,\infty)\) whose derivative is determined by the scaled thresholds.  Specifically, let \(u_k = \exp( \text{something like } k\pi/t )\) after mapping; but a more precise formulation is as follows.

Recall the phase derivative \(\phi'(t)\) should, up to a constant, approximate the zero density.  Define a continuous function \(\psi_G(t)\) by
\[
\psi_G(t) = \sum_{k=0}^{\infty} \alpha_k \delta_k\!\left( \frac{t}{T} \right),
\]
where \(T\) is a large scaling parameter and the indicators are interpreted as functions of \(t\) via a logarithmic change of variable.  After a mollification, we set \(\phi_G'(t) = c\psi_G(t)\) with a constant \(c\) chosen to match the leading asymptotic \(\log(t/2\pi)\).

The precise construction involves:
\begin{enumerate}
\item Map the interval \([0,\pi]\) of \(\theta\) to the logarithmic scale via \(u = \log t\), so that \(\theta = \pi\) corresponds to \(u = \log \lambda\) for some large \(\lambda\).
\item Let \(\delta_n\) act on the scaled coordinate; the thresholds \(\theta_n^*\) become positions \(t_n^*\) whose distribution is governed by the harmonic angles.
\item The derivative \(\phi_G'\) is then a sum of step functions of width and height determined by \(\alpha_n\), which after smoothing yields a smooth phase \(\phi_G\).
\item Compute the torsion \eqref{tors_formula} with \(\phi=\phi_G\) and take the distributional limit as the reconstruction depth \(N\to\infty\) and the mollifier width \(\varepsilon\to0\).
\end{enumerate}

\section{Conjectural Matching}
\begin{conjecture}[Greedy Phase Torsion]
With the above construction, the distributional limit of \(\tors_{\phi_G}\) exists and is given by
\[
\tors_{\phi_G}(t) = \tau_{\text{prime}}(t) + \tau_{\text{arch}}(t).
\]
\end{conjecture}

The intuition is that the singularities of the torsion arise when the derivative \(\phi_G'\) has jumps (from the step functions).  The height of each jump is proportional to an \(\alpha_n\), and the location is at the scaled thresholds.  The correlation kernel \(K(n,m)\) controls the products \(\phi_G''\) and \(\phi_G'''\) that appear in the torsion numerator.  After a careful combinatorial analysis, the kernel‑weighted sum over thresholds should reorganise into the prime sum.

\section{Partial Evidence and Heuristic Calculations}
Preliminary numerical computations for moderate \(N\) (up to a few hundred) indicate that the finite‑\(N\) torsion measure, when integrated against smooth test functions, produces values close to the partial prime sums.  However, a rigorous asymptotic analysis has not been carried out, and the error terms are not yet bounded.

\section{Alternative Approaches}
If the greedy harmonic reconstruction proves too delicate, there may be other ways to prove the Geometric Explicit Formula.  For instance, one could directly analyse the zero‑locking phase \(\phi\) derived from the CCM framework, without recourse to the greedy construction.  The Beurling–Selberg sieve and the theory of prolate spheroidal wave functions might provide the necessary uniform estimates.  The geometric programme is flexible; the essential point is that the conjecture is well‑posed and reduces RH to a clear analytic target.

\chapter{Conclusions and Future Work}

This thesis has presented a geometric research programme that reformulates the Riemann Hypothesis as a conjectural identity between the torsion of a space curve and a prime‑supported distribution.  The main contributions are:

\begin{itemize}
\item Construction of the twin conical helices and the unconditional derivation of the geometric prime distribution \(\tau_{\text{prime}}\) (Chapter 3).
\item Introduction of the Riemann helix and the formulation of the Geometric Explicit Formula (Conjecture~\ref{conj:GEF}), whose truth implies RH (Chapter 4).
\item Clarification that the variable‑radius approach does not yield a proof (Chapter 4).
\item Development of the greedy harmonic sine reconstruction as a candidate for the zero‑locking phase function (Chapter 5).
\item Outline of a programme to prove the Geometric Explicit Formula by scaling the greedy reconstruction to match the zero density (Chapter 6).
\end{itemize}

The central open problem remains:

\begin{openproblem}
Prove the Geometric Explicit Formula (Conjecture~\ref{conj:GEF}), or find a counterexample.
\end{openproblem}

Further research directions include:
\begin{itemize}
\item Rigorous operator construction for the Dirac operator with point interactions in the limit \(\lambda\to\infty\) and the convergence of its spectral determinants.
\item Generalisation to Dirichlet and Hecke \(L\)-functions, with explicit computation of the archimedean background.
\item Numerical exploration of the greedy reconstruction to gather evidence for or against the conjectured matching.
\item Investigation of other combinatorial sieves that might produce a phase function satisfying the torsion identity.
\item Detailed study of the correlation kernel and its relation to the von Mangoldt function.
\end{itemize}

The geometric language offers a fresh perspective on the Riemann Hypothesis, making it accessible to tools from differential geometry, distribution theory, and combinatorics.  Regardless of the ultimate outcome, the programme advanced here provides a valuable new viewpoint and a concrete set of problems for future generations of mathematicians.

\appendix
\chapter{Explicit Computation of the Torsion Denominator}
We include, for completeness, the detailed calculation of the torsion formula \eqref{tors_formula}.

\[
\text{... (full calculation as in previous documents) ...}
\]

\bibliographystyle{plain}
\begin{thebibliography}{9}
\bibitem{CCM}
A.~Connes, C.~Consani, H.~Moscovici,
\emph{The scaling operator and a spectral triple for the Riemann zeta function},
arXiv:2511.22755 (2025).
\bibitem{Albeverio}
S.~Albeverio, F.~Gesztesy, R.~H{\o}egh-Krohn, H.~Holden,
\emph{Solvable Models in Quantum Mechanics}, Springer, 1988.
\bibitem{Weil}
A.~Weil,
\emph{Sur les formules explicites de la th\'eorie des nombres},
Comm. Sem. Math. Univ. Lund, tome suppl. d\'edi\'e \`a Marcel Riesz (1952), 252--265.
\bibitem{Edwards}
H.~M.~Edwards,
\emph{Riemann's Zeta Function}, Academic Press, 1974.
\bibitem{IwaniecKowalski}
H.~Iwaniec, E.~Kowalski,
\emph{Analytic Number Theory}, Amer. Math. Soc., 2004.
\end{thebibliography}

\end{document}